\section*{Capítulo \ref{ch:g5:areas} --- Perímetro y área}

\begin{solution}{exo:g5:areas:1}
\omterm{def:g4:measure:area}{Área}: $8 \times 5 = 40$ cm$^2$. \omterm{def:g4:measure:perimeter}{Perímetro}:
$2 \times (8 + 5) = 26$ cm.
\end{solution}

\begin{solution}{exo:g5:areas:2}
$9 \times 9 = 81$ cm$^2$; \quad $12 \times 7 = 84$ m$^2$; \quad
$4.5 \times 6 = 27$ cm$^2$.
\end{solution}

\begin{solution}{exo:g5:areas:3}
Ancho: $63 \div 9 = 7$ cm. \omterm{def:g4:measure:perimeter}{Perímetro}: $2 \times (9 + 7) = 32$ cm.
\end{solution}

\begin{solution}{exo:g5:areas:4}
Por ejemplo, $6 \times 4$ (\omterm{def:g4:measure:perimeter}{perímetro} $20$) y $8 \times 3$
(\omterm{def:g4:measure:perimeter}{perímetro} $22$): la misma \omterm{def:g4:measure:area}{área} $24$ y \omterm{def:g4:measure:perimeter}{perímetros} distintos.
\end{solution}

\begin{solution}{exo:g5:areas:5}
$3$ m$^2 = 30\,000$ cm$^2$; \quad $50\,000$ cm$^2 = 5$ m$^2$.
\end{solution}

\begin{solution}{exo:g5:areas:6}
Barra: $8 \times 2 = 16$ cm$^2$; pie: $2 \times 5 = 10$ cm$^2$;
total, $26$ cm$^2$.
\end{solution}

\begin{solution}{exo:g5:areas:7}
$30 \times 20 - 24 \times 14 = 600 - 336 = 264$ cm$^2$ de madera.
\end{solution}

\begin{solution}{exo:g5:areas:8}
La \omterm{def:g3:division:half}{mitad} del rectángulo de $8 \times 6$: $48 \div 2 = 24$ cm$^2$.
\end{solution}

\begin{solution}{exo:g5:areas:9}
Semilla: \omterm{def:g4:measure:area}{área} $7 \times 4 = 28$ m$^2$, coste $28 \times 2 = 56$.
Valla: \omterm{def:g4:measure:perimeter}{perímetro} $2 \times (7 + 4) = 22$ m, coste
$22 \times 3 = 66$.
\end{solution}

\begin{solution}{exo:g5:areas:10}
Dos baldosas de $50$ cm hacen un metro: a lo largo de los $12$ m,
$24$ baldosas; a lo ancho de los $2$ m, $4$ baldosas. Total:
$24 \times 4 = 96$ baldosas.
\end{solution}

\begin{solution}{exo:g5:areas:11}
Antes: \omterm{def:g4:measure:perimeter}{perímetro} $2 \times (4 + 3) = 14$ cm y \omterm{def:g4:measure:area}{área} $12$ cm$^2$.
Después (lados $8$ y $6$): \omterm{def:g4:measure:perimeter}{perímetro} $28$ cm y \omterm{def:g4:measure:area}{área} $48$ cm$^2$.
El \omterm{def:g4:measure:perimeter}{perímetro} se ha \emph{duplicado}; el \omterm{def:g4:measure:area}{área} se ha multiplicado por
\emph{cuatro} ($2 \times 2$): las longitudes y las \omterm{def:g4:measure:area}{áreas} no cambian de la misma
manera.
\end{solution}
